% Appendix B: signs and normalisations (label prefix sgn). Conventions and dictionaries only; nothing is proved here.
% Continued in part 2.
\section{Signs and normalisations}
\label{sec:sgn}

This appendix collects in one place the sign conventions and normalisations of the note, and the dictionaries to the conventions of the sources quoted. It proves nothing: every entry is a convention, or a dictionary whose derivation is given at the place cited, and the body refers to it only for orientation.

%-------------------------------------------------------------------------------
\subsection{The \texorpdfstring{$R$}{R}-matrix and its action}
\label{sec:sgn-R}

\emph{Givental's convention.} The $R$-matrix $R(z)=1+O(z)$ of the formal quintic is taken in Givental's convention: the fundamental solution of the quantum differential equation has the form $S\sim R\,\re^{u/z}$, that is, the columns $R(z)\tilde e_i\,\re^{u_i/z}$ solve $z\de_Tv=H*v$ along $t=T(q)H$ (\cref{def:set-rows}), and $R$ is symplectic, $R(-z)^T\eta R(z)=\eta$ in the flat basis. Equivalently $R(z)^{-1}=R(-z)^*$, the adjoint with respect to $\eta$.

\emph{The action.} The action on cohomological field theories is that of \cite[Def.~2.2]{PPZ15}: the legs carry $R^{-1}(\psi)$, and the edge of $R$ is
\[
E_R(\psi,\psi')=\frac{\eta^{-1}-R^{-1}(\psi)\,\eta^{-1}R^{-1}(\psi')^T}{\psi+\psi'},
\]
a power series because $R$ is symplectic \cite[Remark~2.1, (6)]{PPZ15}. The authors of \cite{PPZ15} state after their Def.~2.2 that they use the inverse of the $R$-matrix there and that Givental and Teleman use the same convention. The action of \cite[\S2.3]{Dunin-Barkowski:2015caa} is the same (legs $R^{-1}(\psi)$ acting on column vectors, the row index being the outgoing colour); \cite[Remark~2.1]{Dunin-Barkowski:2015caa} records that Givental's convention for the matrix $R$ is the transpose of theirs. The unit-preserving action is $R.\Omega=R\,\ttr_R\,\Omega$ with the translation
\[
\ttr_R(z)=z\,1-z\,R^{-1}(z)\,1
\]
of \cite[Def.~2.13]{PPZ15}, which is the translation $T_b$ of \cite[Prop.~2.12]{PPZ15}; so the dilaton leaves carry $+\psi\bigl(1-R^{-1}(\psi)1\bigr)$. This is also the sign of the dilaton leaves produced by topological recursion (\cref{prop:mod-Rtr}) and the sign of \cite[\S2.3]{DuninBarkowski:2012bw}; \cite[\S2.3]{Dunin-Barkowski:2015caa} prints the opposite sign. The contraction $\mathrm{Cont}_\Gamma$ includes the factor $|\Aut\Gamma|^{-1}$ (\cref{def:set-graphsum}), so that sums over graphs carry no further factor.

\emph{The colour basis.} Let $\Qr^{(i)}_a(z)=\eta(H^a,R(z)\tilde e_i)$ be the lowered rows (\cref{def:set-rows}), and write the graph sum of \cref{def:set-graphsum} in the basis of idempotents $e_i=\Delta_i^{-1/2}\tilde e_i$, in which $\eta(e_i,e_j)=\delta_{ij}/\Delta_i$. A vertex of genus $g_v$ and colour $i$ carries $\Delta_i^{g_v-1}=\TFT_{g_v,n_v}(e_i,\dots,e_i)$ times the $\psi$-integrals of the vertex; the edge between colours $i$ and $j$ is
\[
E_{ij}(\psi,\psi')=\frac{\delta_{ij}\Delta_i-\Delta_i^{1/2}\Delta_j^{1/2}\sum_{a,b}\eta^{ab}\,\Qr^{(i)}_a(-\psi)\,\Qr^{(j)}_b(-\psi')}{\psi+\psi'},\qquad \eta^{ab}=\tfrac15\,\delta_{a+b\equiv3\,(5)},
\]
and the dilaton leaves are the $e_i$-components of the translation, $\ttr_R(z)=\sum_i\sum_{k\geq2}\mathrm{Dil}_i(k)\,z^k\,e_i$ with $\mathrm{Dil}_i(k)=(-1)^k\Delta_i^{1/2}\coef{z^{k-1}}\Qr^{(i)}_0(z)$. These formulas are \cite[Defs.~2.2, 2.13]{PPZ15} written in the basis $e_i$ (the naturality of \cref{def:set-graphsum} under constant isometries): since $\eta^{-1}=\sum_k\tilde e_k\tensor\tilde e_k$ and $R^{-1}(\psi)=R(-\psi)^*$, the $\tilde e_i\tensor\tilde e_j$-component of $R^{-1}(\psi)\eta^{-1}R^{-1}(\psi')^T$ is $\eta\bigl(R(-\psi)\tilde e_i,R(-\psi')\tilde e_j\bigr)=\sum_{a,b}\eta^{ab}\Qr^{(i)}_a(-\psi)\Qr^{(j)}_b(-\psi')$, and $\tilde e_i\tensor\tilde e_j=\Delta_i^{1/2}\Delta_j^{1/2}\,e_i\tensor e_j$; and the unit is $1=\sum_ie_i$, with $R^{-1}(z)1=\sum_i\Delta_i^{1/2}\Qr^{(i)}_0(-z)\,e_i$ and $\Delta_i^{1/2}\Qr^{(i)}_0(0)=\Delta_i^{1/2}\eta(1,\tilde e_i)=1$. This colour-basis form is the one used in the computations of \cref{sec:evi}; the proofs use the form of \cite{PPZ15}.

\emph{Lowered and vector forms.} Solutions $\phi$ and gauges are written in lowered components $Y_a=\eta(H^a,\phi)$ (\cref{conv:tra-forms}). As a matrix $\eta$ is $5$ times the permutation matrix of $a\leftrightarrow3-a$; a gauge $P$ acting on lowered components acts on vectors by $P^{\mathrm{vec}}=\eta^{-1}P\eta$, the conjugate of $P$ by this permutation, and the dilaton factor is $d:=(P_0)_{33}=(P_0^{\mathrm{vec}})_{00}$ (\cref{def:tra-dressing}). The annihilation condition (A) is stated in lowered form, \cref{lem:tra-E,lem:tra-T} in vector form.

%-------------------------------------------------------------------------------
\subsection{Variables, colours and the hatted theory}
\label{sec:sgn-var}

\emph{Operators.} $\Dq=q\,\dd/\dd q=-(1-u)\,\dd/\dd u$; the derivative in $u$ appears as $z\de_u$, so that $z\Dq=-(1-u)z\de_u$. Along the mirror coordinate $\de_T=C_1^{-1}\Dq$ (\cref{prop:set-cyclicLR}); in \cref{sec:gau-cyc}, $\de_t=(\Dq t)^{-1}\Dq$ (\cref{def:gau-conc}). The Euler operators are $\theta_u=u\,\dd/\dd u$ and $\theta_{\zh}=\zh\,\dd/\dd\zh$ (\cref{lem:gau-HPF}).

\emph{Points.} $u=1-5^5q$, $s=u^{1/5}$, $L=1/s$, $\zh=z/s^4=L^4z$. The large-radius point is $q=0$ ($u=s=L=1$, $\zh=z$), the conifold point is $u=0$, the orbifold point is $q=\infty$ ($L=0$) (\cref{def:set-points}).

\emph{Colours.} Colour $i\in\bbZ/5$ is the fixed point $p_i$ of the torus, of weight $\lambda_i=\zeta^i$; the eigenvalue $\zeta^iL$ of $C_1H*$, so that $\Dq u_i=\zeta^iL$ (\cref{cor:set-cyclicPF5}); and the critical point $\sigma=\zeta^iL$ of the Mellin--Barnes summand (\cref{def:gam-c}). The colour covariances are $\Phi_i(\zh,u)=\Phi_0(\zeta^{-i}\zh,u)$ (\cref{lem:gau-colrows}), $\Dn_i(z)=\Dn_0(\zeta^{-i}z)$ (\cref{cor:gau-colours}) and $h_{\bw_i}(z)=h(\zeta^{-i}z)$ (\cref{lem:set-hrot}). The square roots are fixed by $\zeta^{1/2}:=\zeta^3$: $\Delta_i^{1/2}=\zeta^{4i}C_0/L$, and $(W^0_i)^{-1/2}=\zeta^i$ for $W^0_i=\zeta^{3i}$ (\cref{def:set-frob}).

\emph{Flat and hatted.} The flat theory uses $(u,z)$, the classes $H^a$ and the rows $\Qr$. The hatted theory (\cref{def:set-hat}) uses $(u,\zh)$ with $z=s^4\zh$, the classes $\hat H^a=s^{1-\nu_a}H^a$, the rows $\hat\Qr_a=s^{1-\nu_a}\Qr_a$, the datum $\hat x=\varpi_0$ with $\hat\Delta_i^{1/2}=\zeta^{4i}\hat x$, the unit $\hat1=s\,1$ and the rescaling $\Omega_{g,n}\mapsto s^{-(2g-2+n)}\Omega_{g,n}$; the pairing $\eta$ is preserved, since $\nu_a+\nu_b=2$ on the $\eta$-pairs $(0,3)$, $(1,2)$, $(4,4)$. In the frame $(\alpha,\beta)$, $\hat x=\varpi_0=1+\alpha u+O(u^2)$. The hatted term of a leg-free graph of genus $g$ is $G_\Gamma=u^{2g-2}T_\Gamma$ (\cref{lem:flh-hat}).

%-------------------------------------------------------------------------------
\subsection{Quantum Riemann--Roch}
\label{sec:sgn-qrr}

\emph{The twist.} In the weights of quantum Riemann--Roch at the fixed point $p_i$ the tangent directions contribute their weights $\lambda_i-\lambda_j$, $j\neq i$, and an Euler-class twist contributes minus its fibre weight; the fibre weight of $\cO(5)$ at $p_i$ is $5\lambda_i$ \cite[(10)]{LP20}. So at $\lambda_i=\zeta^i$ the weights of (QRR) are $\bw_i=\{\zeta^i-\zeta^j\}_{j\neq i}\cup\{-5\zeta^i\}$, with $h_{\bw}(z)=\exp\bigl(-\sum_k\bk_kz^{2k-1}\sum_{a\in\bw}a^{1-2k}\bigr)$ and $\bk_k=B_{2k}/(2k(2k-1))$ (\cref{thm:set-qrr}). Since $\sum_{a\in\bw_0}a^{-1}=\sum_{j=1}^4(1-\zeta^j)^{-1}-\frac15=\frac95$, the first coefficient is $h_1=-\frac3{20}$; the weight $+5\zeta^i$ would give $-\frac{11}{60}$ (\cref{rmk:set-qrrsign}). Quantum Riemann--Roch is used in the form $\Phi_i(z,1)=h_{\bw_i}(z)$ (\cref{eq:set-qrr}).

\emph{Characters.} \cite[\S2.2, Theorem~1]{Giv01} is written with the tangent characters at the fixed points and a minus sign in the exponent, \cite[9.1~Theorem]{MR1901075} with the cotangent characters and a plus sign; the two agree, since only odd powers occur.

\emph{Bernoulli numbers.} The Bernoulli numbers and polynomials are those of \cref{def:fsp-stirling}: $\xi/(\re^\xi-1)=\sum_nB_n\xi^n/n!$ and $\xi\,\re^{\theta\xi}/(\re^\xi-1)=\sum_nB_n(\theta)\xi^n/n!$ \cite[24.2.3]{DLMF}, so $B_n=B_n(0)$. Of the Bernoulli numbers only the $B_{2k}$, $k\geq1$, enter the note, so the conventions $\xi/(\re^\xi-1)$ and $\xi/(1-\re^{-\xi})$ of the literature, which differ only in $B_1$, give the same formulas; the polynomials $B_n(\theta)$ are always those of \cite[24.2.3]{DLMF}.

\emph{Normalisation of asymptotic series.} \cite[Theorem~4 (i)]{ZZ08} and \cite[\S2.4.4, (18)]{LP20} normalise the asymptotic series attached to the $I$-function so that they equal $1$ at $q=0$ ($\Phi_s(0)=\delta_{0,s}$ there; in \cite[\S2.4.4]{LP20} the first coefficients of (18) are printed as $R_0=L$, $R_1=\frac3{20}(L-L^5)$, $R_2=\frac{9L}{800}(1-L^4)^2$, with the values $1,0,0$ at $L=1$). This is not (QRR), which gives the value $h_{\bw_i}(z)$ at $q=0$.

%-------------------------------------------------------------------------------
\subsection{Normalisations of the model}
\label{sec:sgn-model}

The note uses one normalisation of the model, $\Wv(p)=p^4/4+v/p$ with the integrands $p^m\re^{-\Wv(p)/z}\dd p$ and the curve $x=\Wv(p)$, $y=p$ (\cref{def:set-model,def:mod-curve}); its branch points are $\varsigma_b=\zeta^{-b}\varsigma$, $\varsigma^5=v$, with $x''(\varsigma_b)=5\varsigma_b^2$ and the pairwise distinct critical values $x(\varsigma_b)=\frac54\zeta^b\varsigma^4$. The scaling $p\mapsto kp$, $v\mapsto k^5v$, $z\mapsto k^4z$ is a symmetry of the model. \Cref{tab:sgn-model} gives the other normalisations that occur in the sources quoted.

\begin{table}[ht]
\caption{Normalisations of the model.}
\label{tab:sgn-model}
\centering
\footnotesize
\renewcommand{\arraystretch}{1.2}
\begin{tabular}{@{}>{\raggedright\arraybackslash}p{0.24\textwidth}>{\raggedright\arraybackslash}p{0.4\textwidth}>{\raggedright\arraybackslash}p{0.3\textwidth}@{}}
\hline
form & relation to $\Wv$ & use \\
\hline
\cite[(8-1)]{Dunin-Barkowski:2015caa}, $n=4$: $f=p^4+a_2p^2+a_3p+a_4+\frac{a_5}{p-a_1}$ & at $a=(0,0,0,0,4v)$, $f=4\Wv$; critical points $p^5=v$, critical values $5p_c^4$ & the quoted theorem of \cref{thm:mod-tr}, transferred to $x=\Wv$ by \cref{lem:mod-scale} \\
the curve $x=p^4+V/p$, $y=p$ (cf.\ \cite[(7.2)]{ABDKS25}) & $V=4v$: $x=4\Wv$; $\omega_{h,n}[\Wv]=4^{2h-2+n}\,\omega_{h,n}[p^4+V/p]|_{V=4v}$, and every pairing scales by $4^{2h-2+n}$ & not used: \cite[Theorem~7.1]{ABDKS25} is applied to $x=\Wv$ itself (\cref{lem:mod-abdks}) \\
$\lambda=x^4+u'/x$, integrand $\re^{\lambda/z'}x^k\dd x$ & $x=p$, $u'=4v$, $z'=-4z$: $\re^{\lambda/z'}=\re^{-\Wv/z}$; at $v=1$ the row $b$ of colour $i$ of \cref{def:set-M0} is $\varepsilon_b\zeta^{i(b+1)}J_{\nu_b}(-4\zeta^{-i}t)$, where $J_k(z')$ (a letter of this row) is the normalised stationary-phase series $\SPn_{x=1}[\lambda;x^k](z')$ of this integrand at its critical point $x=1$ (\cref{prop:set-modrows,lem:con-S2}) & dictionary only \\
$\lambda=-\Wv$, integrand $\re^{\lambda/z}$ & the same integrand & dictionary only \\
TFT of topological recursion & $\Delta_j=-5\varsigma_j^2$ in the basis $e^{\mathrm{TR}}_j$ of \cite[(2-30)]{Dunin-Barkowski:2015caa}, with $x=x(\varsigma_j)-w_j^2/2$ and the uniform orientation $\frac{\dd w_j}{\dd p}(\varsigma_j)=\sqrt{-5}\,\varsigma_j$ (one root for all $j$); frozen basis $\bar e_j=e^{\mathrm{TR}}_j/\sqrt{-5}$ & \cref{prop:mod-Rtr,thm:mod-tr} \\
model TFT & defined by $\hat\Delta_i^{1/2}:=\zeta^{4i}\varsigma$ in the normalised canonical basis, unit $\sum_i\hat\Delta_i^{-1/2}\tilde e_i$ & \cref{def:set-modcohft} \\
topological recursion versus the model & $\omega_{h,n}=(-5)^{h-1}\sum\langle\prod_i\psi^{d_i}H^{a_i}\rangle_h\prod_i\dd\xi_{a_i,d_i}$; $\Fm_h=(-5)^{1-h}\FTR_h$ & \cref{thm:mod-tr,cor:mod-trF} \\
Hurwitz space & \cite[(8-1)]{Dunin-Barkowski:2015caa} with $n=4$ on the line $a=(0,0,0,0,4v)$; called $H_{0;1,4}$ in \cite{CvdLPS23} & the identification is not used \\
$H$ and $\se$ & $H=\dd/\dd v$, $\se=z\,1+vH=\fr(1,0)$; $\langle H\rangle_1=-1/(12v)$, $\langle\se\rangle_1=\frac18$; replacing $H$ and $\se$ by $-H$ and $-\se$ changes both signs & \cref{lem:tra-frames}; sign anchors, used in no proof \\
\hline
\end{tabular}
\end{table}

% Appendix B, part 2: topological recursion, legs, the tower, normalisations, stationary phase, the literature.

%-------------------------------------------------------------------------------
\subsection{Topological recursion}
\label{sec:sgn-tr}

Topological recursion is used in \cref{sec:mod} only, on the curve $x=\Wv(p)$, $y=p$, with the kernel of \cref{conv:mod-kernel}; $\iota_b$ is the local involution of $x$ at $\varsigma_b$. \Cref{tab:sgn-tr} compares it with the printed conventions of the sources quoted.

\begin{table}[ht]
\caption{Conventions of topological recursion.}
\label{tab:sgn-tr}
\centering
\footnotesize
\renewcommand{\arraystretch}{1.2}
\begin{tabular}{@{}>{\raggedright\arraybackslash}p{0.15\textwidth}>{\raggedright\arraybackslash}p{0.4\textwidth}>{\raggedright\arraybackslash}p{0.39\textwidth}@{}}
\hline
object & this note (\cref{conv:mod-kernel}) & sources \\
\hline
kernel & $K_b(p_0,q)=\frac12\int_{\iota_bq}^{q}\frac{\omega_{0,2}(p_0,\cdot)}{(y(q)-y(\iota_bq))\,\dd x(q)}$ near $\varsigma_b$, with $\omega_{0,1}=y\,\dd x$ and $\omega_{0,2}=\frac{\dd p_1\dd p_2}{(p_1-p_2)^2}$ & equal to the kernel of \cite[(2-26)--(2-29)]{Dunin-Barkowski:2015caa} and of \cite[(2.1)]{Hock25}; the printed kernel of \cite[(4-3), (4-9)]{Eynard:2007kz} equals that of \cite[(1.1)]{ABDKS25} with $\omega^{(0)}_{1,0}=-y\,\dd x$, and both are minus this kernel \\
correlators & $\omega_{h,n}$ & $\omega_{h,n}=(-1)^nW^{\mathrm{EO}}_{h,n}=(-1)^n\omega^{\mathrm{ABDKS}}_{h,n}$ \\
$\omega_{0,3}$ & $-\sum_b\Res_{q=\varsigma_b}\frac{\omega_{0,2}(q,p_1)\,\omega_{0,2}(q,p_2)\,\omega_{0,2}(q,p_3)}{\dd x(q)\,\dd y(q)}$ & minus \cite[Theorem~4.1, (4-20)]{Eynard:2007kz} \\
free energies & $(2-2h)\FTR_h=\sum_b\Res_{\varsigma_b}\Phi_{y\mathrm dx}\,\omega_{h,1}$, $\dd\Phi_{y\mathrm dx}=y\,\dd x$; on the model curve $\FTR_h=\frac{B_{2h}}{2h(2h-2)}v^{2-2h}$ for $h\geq2$ (\cref{thm:mod-gap}) & $\FTR_h=-F^{\mathrm{print}}_h$, with $F^{\mathrm{print}}_h$ of \cite[(4-10)]{Eynard:2007kz} \\
dilaton equation & $\sum_b\Res_{q=\varsigma_b}\Phi_{y\mathrm dx}(q)\,\omega_{h,n+1}(q,p_1,\dots,p_n)=(2-2h-n)\,\omega_{h,n}(p_1,\dots,p_n)$ & \cite[Theorem~4.7, (4-30)]{Eynard:2007kz}, with the factor $2h+n-2$ for $W^{\mathrm{EO}}$ \\
variation in $v$ & \cref{thm:mod-var}, at fixed $x$, along a ray $\mathfrak r$ & \cite[Theorem~5.1, (5-14)--(5-15)]{Eynard:2007kz} with \S5.3, (5-27)--(5-30); for $\partial_v$ at fixed $x$ the function $\Lambda$ of \cite[(5-13)]{Eynard:2007kz} (a letter of that paper) is $-1$ on the ray, with the printed signs of \cite[(5-1), (5-13)]{Eynard:2007kz} \\
Laplace transforms & \cite[(2-31), (2-32)]{Dunin-Barkowski:2015caa} read as formal Laplace series, the finite part at $i=j$ giving the $\delta_{ij}$ of $R^{-1}=1+O(\vartheta^{-1})$ & \cref{lem:mod-laplace} \\
sign anchors & the curve $x=-p^2/2$, $y=p$ gives the intersection numbers of $\psi$ classes; $x=p+1/p$, $y=p$ gives $\FTR_2=-\frac1{240}$ (\cref{conv:mod-kernel}) & classical; used in no proof \\
\hline
\end{tabular}
\end{table}

%-------------------------------------------------------------------------------
\subsection{Legs}
\label{sec:sgn-legs}

The dictionary of record is the leg dictionary of \cref{prop:mod-LD}. With the legs $\phi^b_0=p,\log p,v/p,p^3,p^2$ ($b=0,\dots,4$) and $\phi^b_{k+1}=-\int\phi^b_k\,\dd x$ (\cref{def:mod-legs}), one has $\langle\phi^b_k,\dd\xi_{a,d}\rangle=-\varrho_b\delta_{ab}\delta_{dk}$ with $\varrho=(-1,1,-1,-3,-2)$; hence $q_{\phi^b_k}=-\varrho_b\,\psi^kH^b$, and the leg of $v^m\psi^kH^b$ is $-\varrho_b^{-1}v^m\phi^b_k$. In particular the unit has the leg $y=p$, and in the string equation it acts on the legs by $+\dd/\dd x$ (an orientation; the string equation is used in no proof); the dilaton shift $q=-z\,1$ has the leg $\int y\,\dd x=-\phi^0_1$; the class $H$ has the logarithmic leg $-\log p$; and the Euler--dilaton vector $\se=z\,1+vH$ has the leg $-p^5/5$.

Legs in the leg class $\cL=\bbQ[v][p]+\bbQ(v)[x]$ are paired with $\omega_{h,n}$ by the residue pairing $\sum_b\Res_{\varsigma_b}$ (\cref{def:mod-pair}). For legs in $\bbQ[v][p]$ it equals the pairing on a circle $|p|=R>|\varsigma|$; for legs with a pole at $p=0$ the two differ by the residue at $p=0$, which is not zero in general. Logarithmic legs are paired along a ray $\mathfrak r$ from $0$ to $\infty$ that avoids the branch points.

A second dictionary, with the unit leg $-p$, the leg $-\log p$ for $H$ and the step $+\int\cdot\,\dd x$, is not dual to the functions $\xi_{a,d}$ and is not used; on the curve $x=p^4+V/p$ it gives $\se$ the leg $-\frac45p^5$ and $\psi\se$ the leg $-\frac{16}{45}p^9+\frac V5p^4$.

%-------------------------------------------------------------------------------
\subsection{The vanishing tower and regularity}
\label{sec:sgn-tower}

The note uses one tower of covectors, defined from the vanishing solution $\YC$ of the model (\cref{def:tra-tower,def:gau-van}):
\[
\Pi_i(q):=-\coef{z^i}\bigl\langle\YC(-z),q(z)\bigr\rangle,\qquad i\geq0;
\] $\Pi_0$ is the vanishing covector, and $\Pi_i(z^kH^b)=\Pi_0(z^{k-i}H^b)$, which is zero for $k<i$. On insertions, $\Pi_0(z^kH^b)=(-1)^{k+1}\coef{z^{-k}}\YC_b$, which vanishes unless $b+k\equiv1\bmod5$; for instance $\Pi_0(H)=1$ and $\Pi_0(z\,1)=-v$.

\emph{Weights.} The $v$-exponent of an ancestor correlator $\langle\prod_i\psi^{k_i}H^{a_i}\rangle_h$ of the model is $e=2-2h+\sum_i\wt(k_i,a_i)$ with $\wt(k,a)=(4k+\nu_a-4)/5$ (\cref{lem:set-modhom}). The same exponent is also written $2-2h-\sum_i(4-\nu_{a_i}-4k_i)/5$, where $(4-\nu_a)_{a=0,\dots,4}=(4,5,1,2,3)$. In the graded setting of \cref{conv:tra-forms} the variable $v$ has weight $5$ and $z^kH^a$ has weight $4k+\nu_a+1$, so that the weight of $z^kH^a$ is $5+5\wt(k,a)$.

\emph{Regularity.} A correlator of the model is regular if its $v$-exponent is $\geq0$, that is, if it lies in $\bbQ[v]$; a graph term is regular if it lies in $\bbQ\fps{u}$; for the hatted term $G_\Gamma=u^{2g-2}T_\Gamma$ of a genus-$g$ graph, regularity of $T_\Gamma$ means $G_\Gamma\in u^{2g-2}\bbQ\fps{u}$. ``Regular'' never means ``zero''.

%-------------------------------------------------------------------------------
\subsection{Normalisations of gauges and series}
\label{sec:sgn-norm}

The normalisations used throughout are: the canonical gauge has $P_{11}(0,z)=1$ (\cref{def:gau-can}); the vanishing period is $\varpi_1=u+O(u^2)$, and $\varpi_0(0)=1$, $\tc'(0)=1$ in the conifold frames (\cref{lem:set-frobbasis,def:set-frames}); $v(u)=u+O(u^2)$ (\cref{def:set-tau}); $\Fs_0(0)=1$ (\cref{prop:set-modrows}); $\mathrm{ph}_0=1$ (\cref{def:set-rows}); $c_0(L)=1$ (\cref{def:gam-c}); the series $\SPh$ of \cref{def:fsp-wick} have constant term $1$ (the series $\SPn$ has the constant term $a_0(c)$, \cref{sec:sgn-sp}); and the dilaton factor of $\Pc$ satisfies $d(0)=1$ (\cref{lem:gau-diag}).

The other normalisation of the gauge, used only in \cref{sec:evi}, is fixed by $p_{5n,0}(0)=0$ for $n\geq1$, where $P_k=\sum_bp_{k,b}E_{b+k,b}$ (\cref{def:gau-onepole}); it equals $\Pc\,\cn(z)$ with $\cn(z)=1+m_5z^5+m_{10}z^{10}+\dots$ (\cref{rmk:mai-stored}), and the rational numbers $m_5$, $m_{10}$ are recorded in \cref{sec:evi-univ}.

%-------------------------------------------------------------------------------
\subsection{Stationary phase}
\label{sec:sgn-sp}

\emph{Phases.} Stationary-phase integrands are written $\re^{F/z}$ in \cref{sec:fsp}; the saddle series of \cref{def:gam-c} has the phase $F=\mathsf S_0$; the flat integrand of \cref{sec:con} is written $\re^{-\cW/z}$ and the hatted one $\re^{-\hat\cW/\zh}$ (\cref{lem:con-form}, \cref{def:con-PhiSP}). They correspond under $F=-\cW$, resp.\ $F=-\hat\cW$ with $z$ renamed $\zh$ (\cref{def:fsp-wick}).

\emph{Measures.} In the Mellin--Barnes summands (\cref{sec:con-form}) the Mellin variable is $\sigma$. $\Gamma(1+5\sigma/z)=\int y^{5\sigma/z}\re^{-y}\dd y$ carries the measure $\dd y$, and $\Gamma(w)=\int x^w\re^{-x}\dd x/x$ the measure $\dd x/x$. An insertion $x^m$ with the measure $\dd x/x$ is the case $M=m$ of \cref{lem:fsp-stirling}, with the measure $\dd x$ the case $M=m+1$.

\emph{Normalised and full stationary phase.} The series $\SPn$ of \cref{def:fsp-wick} is normalised by the Gaussian, so that its constant term is the value $a_0(c)$ at the critical point of the $z^0$-part $a_0$ of the amplitude; $\SPh=\SPn/a_0(c)$ has constant term $1$; the full symbol is $(2\pi z)^{n/2}\det(-\mathsf H)^{-1/2}\re^{F(c)/z}\SPn$, with $\mathsf H$ the Hessian at $c$. Only normalised series are canonical: a $z$-dependent prefactor such as $(2\pi z/5)^{1/2}$ is common to the two sides of a comparison and cancels (\cref{def:con-PhiSP}).

\emph{Branches.} Roots and logarithms occur only in the symbols: $\det(-J^{\mathrm T}\mathsf HJ)^{1/2}:=\det J\cdot\det(-\mathsf H)^{1/2}$ under a change of variables with Jacobian $J$ (\cref{lem:fsp-cov}), $\det(-\mu\mathsf H)^{1/2}:=\mu^{n/2}\det(-\mathsf H)^{1/2}$ under a rescaling by $\mu$ (\cref{lem:fsp-homog}), $\log(a/z):=\log a-\log z$ and $\Gst(1+w):=w\,\Gst(w)$ (\cref{def:fsp-stirling}). The saddle series $c(z;L)$ does not depend on the branches (\cref{prop:gam-rational}).

\emph{$q$ and $-q$.} In every form of the Mellin--Barnes summand the critical equation is $-q(5\sigma)^5=\prod_j(\zeta^j-\sigma)$, that is $u\sigma^5=1$, so that $\sigma=L$ is the critical point of colour $0$; the logarithm $\lambda$ of $-q$ is forced by criticality (\cref{def:gam-c}). The series $c(z;L)$ depends on neither $\lambda$ nor the sign convention for $q$, which only fix the dictionary $q=(1-L^{-5})/5^5$. The two summands $(-q)^{\sigma/z}\Gamma(1+5\sigma/z)\prod_j\Gamma((\zeta^j-\sigma)/z)$ and $q^{\sigma/z}\Gamma(1+5\sigma/z)/\prod_j\Gamma(1+(\sigma-\zeta^j)/z)$ give the same normalised series (\cref{lem:con-form}). In Theorem~LMI the chart parameter $\mathsf Q=\re^{-\lambda/a_*}$ is $-q$ at $s_*=1$ and $(-q)^{-1/5}$ at $s_*=0$ (\cref{def:gam-lmidata,thm:gam-main,thm:gam-ginf}).

\emph{Exponentials.} In the comparison of the actual rows with the canonical gauge the same exponential $\re^{u_i/z}$ is removed on both sides (\cref{thm:gau-R0}). The hatted phase gives $\re^{-\hat\cW(1)/\zh}=\re^{(u_0-u_0^{\mathrm c})/z}$, which differs from $\re^{u_0/z}$ by the constant factor $\re^{-u_0^{\mathrm c}/z}$, not a power series in $z$; for this reason the normalisations of \cref{sec:con} are expressed through normalised series and never through the exponential.

%-------------------------------------------------------------------------------
\subsection{Conventions of the literature}
\label{sec:sgn-lit}

\Cref{tab:sgn-lit} records the conventions of the sources in their own letters; in the note they are used only through the translations given at the places cited.

\begin{table}[ht]
\caption{Conventions of the literature and of the note.}
\label{tab:sgn-lit}
\centering
\footnotesize
\renewcommand{\arraystretch}{1.2}
\begin{tabular}{@{}>{\raggedright\arraybackslash}p{0.14\textwidth}>{\raggedright\arraybackslash}p{0.47\textwidth}>{\raggedright\arraybackslash}p{0.33\textwidth}@{}}
\hline
topic & literature & this note \\
\hline
$R$-matrix action & \cite[Def.~2.2]{PPZ15}, the convention of Givental and Teleman; \cite[\S2.3]{Dunin-Barkowski:2015caa} the same action, with Givental's matrix the transpose of theirs \cite[Remark~2.1]{Dunin-Barkowski:2015caa}; the dilaton leaves of \cite[\S2.3]{Dunin-Barkowski:2015caa} have the sign opposite to \cite[Def.~2.13]{PPZ15} and \cite[\S2.3]{DuninBarkowski:2012bw} & \cite{PPZ15} throughout (\cref{sec:sgn-R}) \\
characters & \cite[\S2.2, Theorem~1]{Giv01}: tangent characters, minus sign; \cite[9.1~Theorem]{MR1901075}: cotangent characters, plus sign & equal (\cref{sec:sgn-qrr}) \\
twist weight & \cite[(17)]{GJR18}: $-5h_\alpha$; \cite[(20)]{GJR17}: $t-5H$; \cite[Corollary~1]{MR2276766}: the symbol $\exp\bigl(+\sum_m\bk_m(z/(\lambda+\rho))^{2m-1}\bigr)$ of an Euler-class twist; \cite[(26)]{LP20}: $(5\lambda_i)/e(\mathbb E^*\tensor(5\lambda_i))$ & all agree with $-5\zeta^i\in\bw_i$ (\cref{thm:set-qrr,rmk:set-qrrsign}) \\
asymptotic series & \cite[Theorem~4 (i)]{ZZ08}, \cite[\S2.4.4, (18)]{LP20}: $\Phi_s(0)=\delta_{0,s}$ & (QRR): the value $h_{\bw_i}$ at $q=0$ \\
Landau--Ginzburg charts & \cite[\S7, Remark~7.2]{GJR18}: critical point $(x_i)_\alpha=L_\alpha-\lambda_i$, $y_\alpha=mL_\alpha$, measure $\dd x/x$ with $x^{k/m}$; \cite[\S10]{MR1901075}: chart $x_0=Q/(x_1\cdots x_n)$, measure $\dd x$; \cite[Lemma~4.20]{MR2553377}: $\re^{-W/z}$ at $z=-1$, measure $\dd w/w$; \cite[\S7, (11)]{MR2276766}: $\dd x$ with $(2\pi z(\lambda+\rho))^{-1/2}$ & the seven-variable chart (\cref{rmk:gam-sevenvar}): critical point $x_j=\zeta^j-L$ (the opposite sign to $L_\alpha-\lambda_i$) and $y=5L$, with $\dd y$ for the factor $\Gamma(1+5\sigma/z)$ and $\dd x_j/x_j$ for the factors $\Gamma((\zeta^j-\sigma)/z)$ \\
gap constants & \cite[(1.2)]{HKQ09}: $\frac{(-1)^{g-1}B_{2g}}{2g(2g-2)}t_D^{2-2g}$; \cite[Conjecture~1.5]{CGYZ26} (quintic): $\frac{(-1)^{g-1}B_{2g}}{5^{g-1}2g(2g-2)}\tau_c^{2-2g}$; \cite[Definition~1.8]{CGYZ26} has $(-1)^g$ with a constant $\kappa$ and is not used; \cite[(1.5)]{Br25} (local $\bbP^2$): $\frac{3^{g-1}B_{2g}}{2g(2g-2)}t_{\mathrm{CF}}^{2-2g}$ & $c_g=\frac{(-1)^{g-1}B_{2g}}{2g(2g-2)5^{g-1}}$ for $\varpi_1=u+O(u^2)$; rescaling $\varpi_1$ by $\sqrt5$ gives the form of \cite[(1.2)]{HKQ09} (\cref{rmk:mai-constant}); $\Fm_h=(-5)^{1-h}\FTR_h=c_hv^{2-2h}$ \\
mirror theorem & \cite[Corollary~11.7]{MR1408320} prints $l_\alpha$ for $\lambda_\alpha$ & \cite[Theorem~11.6]{MR1408320} is used (\cref{thm:set-mirror}) \\
\hline
\end{tabular}
\end{table}

